\section{Introduction}
A critical self-avoiding walk crossing a strip of height $h$ is expected
to have length of order $h^{4/3}$. A comparison of total mass and first
length moment suggests this scale, but cannot locate a substantial mass
of paths there: a small collection of very long paths could carry the
entire first moment. This paper supplies a quantitative window by combining
an exactly normalized cylinder calculation with a positive second-moment
and sewing argument.

We realize the honeycomb lattice as the centers of the unit equilateral
triangles in the triangular tiling. A \emph{port} is the midpoint of a
triangle side. A path between ports visits a self-avoiding sequence of
triangle centers and has critical weight
\[
 \mu(\omega)=X^{|\omega|},\qquad X=(2+\sqrt2)^{-1/2},
\]
where $|\omega|$ counts the visited centers. Gluing two paths at a port
multiplies their weights. Lengths of vertex-rooted walks counted by edges
have a different endpoint convention; all statements here use port paths
and center length.

Fix one lattice-side direction, and number its parallel grid lines by
integer height, with one interline distance as unit. A \emph{strict bridge}
of height $H\ge1$ starts at a fixed port on line $0$, ends at any port on
line $H$, and has all centers strictly between the two lines. Let
$\mathcal B_H$ denote these bridges. The start is fixed separately in
each sum; the terminal lateral position is free. Translations and lattice
symmetries identify the possible starts. Write
\[
 B_H=\sum_{\omega\in\mathcal B_H}\mu(\omega),\qquad B_0=1.
\]

\Needspace{12\baselineskip}
\begin{theorem}[A logarithmic bridge-length window]\label{L:thm:window}
There are constants $c,C>0$ and $h_0$ such that, for every real
$h\ge h_0$, with
\[
 R_h=h(\log h)^{1/64},\qquad
 \ell_-(h)=h^{4/3}(\log h)^{-1/8},\qquad
 \ell_+(h)=h^{4/3}(\log h)^{1/2},
\]
we have
\[
 \sum_{1\le H\le 2R_h}
 \sum_{\substack{\omega\in\mathcal B_H\\
       \ell_-(h)\le|\omega|\le2\ell_+(h)}}\mu(\omega)
 \ge c h^{3/4}(\log h)^{-C}.
\]
The sum over $H$ is over integer lattice heights.
\end{theorem}

The averaging in height is part of the conclusion. A particular height,
lateral endpoint, or exact center length is not selected by this theorem.
The window is useful for subsequent renewal conditioning because its length
scale is explicit and its loss is only logarithmic. Passing from such a
window to one prescribed length requires a further local renewal argument.

Several finite estimates have independent uses. For an arch in a half-plane,
let $A_g$ be its critical mass from one fixed boundary port to the boundary
port at positive gap $g$, in lattice-side units. The path has no intermediate
boundary contact. For a plane polygon $P$, length again counts its vertices;
unrooted sums below are modulo primitive lattice translations.

\begin{theorem}[Finite boundary and polygon estimates]\label{L:thm:finite}
There are positive constants, independent of the integer parameters, such
that
\[
 B_H\asymp H^{-1/4},\qquad A_g\asymp g^{-5/4}\quad(H,g\ge1).
\]
The lower bound for $A_g$ can be restricted to arches of diameter at most
$C g$, for a fixed sufficiently large $C$.
For every sufficiently large fixed $T$, bridges of height $H$ whose
transverse displacement from the starting normal stays at most $TH$
have mass at least $c_T H^{-1/4}$ and length-weighted mass at most
$C_T H^{13/12}$. Either weak sign of terminal transverse displacement
may be specified in the mass lower bound.
For arches and bridges from a fixed wall port in a strip of height $H$,
the total mass of paths with diameter exceeding $r\ge C H$ is at most
\[
 C H^{-1/4}e^{-c r/H}.
\]
Finally, for all sufficiently large $R$,
\[
 \sum_{R\le\diam P<2R}\mu(P)|P|^2
 \le C R^{2/3}(1+\log\log R).
\]
\end{theorem}

Here and throughout, $f\asymp g$ means two inequalities with positive
constants independent of the size parameter. Constants may depend on
explicitly fixed confinement or angular parameters. The first-moment
bound in Theorem~\ref{L:thm:finite} has fixed transverse confinement; later
height-averaged moment bounds allow larger diameter cutoffs with their
stated losses.

\subsection{The analytic obstacle and the geometric transfer}
The computation starts with finite loop diagrams on a cylinder. Pairing
states at a horizontal cut give a finite transfer matrix, and diagram
slides determine its vacuum vectors by polynomial interpolation. Two
scalar polynomials, a Pfaffian $Q$ and a determinant $R$, encode the
normalizations. Their definitions and the exact physical gluing formulas
appear in Section~\ref{L:part1}.

A formal asymptotic exponent would be insufficient: the coefficient of
its leading term could vanish. Sections~\ref{L:part2}--\ref{L:part4}
address this issue directly. An oscillator trace is evaluated in two
Gaussian channels; coercivity of the short-range interactions justifies
the change of expansion. A long-circle argument produces expansions with
controlled remainders. Positive Fourier measures then express confluent
$Q$ and $R$ in terms of polynomial norms and force their leading
coefficients to be nonzero. One physical consequence is that the mass of
winding polygons, counted modulo row translation on the regular
$n$-column cylinder, satisfies
\[
 J_n\sim\frac{\sqrt3}{6n}.
\]
Section~\ref{L:part5} strengthens this information to an exponential tail
in the long direction. Section~\ref{L:part6} gives a separate open-strip
calculation of $B_H$; its boundary vectors and determinant are proved in
their own geometry.

The next question is how cylinder observables count long planar paths.
An enclosing-loop gas with loop fugacity $2$ converts, by a local flux
reversal, to the mass of boundary chords visiting a given vertex
(Section~\ref{L:part7}). A two-port identity
(Section~\ref{L:part8}) counts two points on one polygon. Its oblique
asymptotic in Section~\ref{L:part9} is uniform even when the two intervals
between the marks have different scales. The exponent controlling the
scale ratio is independent of the requested expansion accuracy. This
uniformity permits calibration at a smaller physical scale and yields
the polygon second-moment estimate in Theorem~\ref{L:thm:finite}.

Sections~\ref{L:part10} and \ref{L:part11} complete the positive geometric
argument. Strip crossing mass supplies confined paths and caps at
specified gaps. The enclosing gas gives a lower first moment, while the
polygon second moment controls the loss from pruning long paths. Paths
of excessive aspect ratio are excluded by the cylinder pressure tail.
The surviving paths are either bridges already, or arches to which we
attach two disjoint exterior connectors. A bound on the number of double
cuts controls recovery multiplicity. Summing the resulting bridges over
the stated range of heights proves Theorem~\ref{L:thm:window}.

\subsection{Earlier work and the role of exact normalization}
Nienhuis's analysis of the dilute $O(n)$ model predicts the
self-avoiding-walk spatial exponent $3/4$ at $n=0$
\cite{Nienhuis1982}. Lawler, Schramm and Werner formulated the
$\mathrm{SLE}_{8/3}$ scaling-limit and boundary-weight $5/8$ predictions
\cite[Sections~3.3--3.4 and~4.1]{LawlerSchrammWerner2004}.
Duminil-Copin and Smirnov proved the exact honeycomb connective constant
using a parafermionic observable whose local cancellation becomes a
positive boundary identity after summation
\cite{DuminilCopinSmirnov2012}. They also recorded the predicted boundary
powers $A_g\asymp g^{-5/4}$ and $B_H\asymp H^{-1/4}$ in their discussion
of conformal invariance \cite[Section~4]{DuminilCopinSmirnov2012}.
The present finite estimates establish these bounded-factor powers in
the port conventions above; the height-averaged window then locates
critical mass near the corresponding length scale.

The decay of the strip crossing mass has been studied directly.
Beaton, Bousquet-M\'elou, de Gier, Duminil-Copin and Guttmann proved
$B_H\to0$ in their study of surface adsorption
\cite[Theorem~10]{BeatonEtAl2014}. Glazman and Manolescu gave a shorter
proof and a logarithmic subsequence bound
\cite[Proposition~1.1]{GlazmanManolescu2017}. Krachun and Panagiotis
proved a polynomial upper bound for $B_H$ and quantitative
sub-ballisticity for uniform honeycomb walks
\cite[Theorems~1--2]{KrachunPanagiotis2026}. Their use of bridge renewal
levels to recover sewn walks is also a precedent for the positive
geometric counting below. Our constructions keep the cuts recoverable
while imposing the length and confinement bounds needed for the window.

The rhombic weights and their Yang--Baxter relations belong to the
integrable dilute-loop setting. Ikhlef and Cardy related the local
parafermionic equations to the integrable weights
\cite[Section~3]{IkhlefCardy2009}; Glazman established their criticality
for self-avoiding walks on rhombic tilings
\cite[Theorems~1.1--1.2]{Glazman2015}. Glazman and Manolescu proved
boundary two-point invariance for columnwise rhombic half-plane tilings
with angles in $[\pi/3,2\pi/3]$
\cite[Theorem~1]{GlazmanManolescu2017}. The periodic spectral and seam
formulation has precedents in the dilute-model analysis of Zhou and
Batchelor \cite[Section~2]{ZhouBatchelor1997}. We derive the finite
identities in our normalization and then prove the nonzero amplitudes,
uniform errors and planar conversions used here.

The oscillator fields, cocycles and root-mode identities in the change
of expansion come from lattice vertex-operator algebra
\cite[Section~2.1]{Frenkel1985}. For the present indefinite bilinear
space, we establish finite coefficient support before taking traces,
then prove absolute summability and continuation. The pressure argument
also uses classical total positivity and variation diminution
\cite{Schoenberg1947}, followed by the connected-cluster sign structure
of a repulsive gas \cite[Sections~2.2--2.3]{ScottSokal2005}.
The local proofs supply the strictness and bounds uniform in the
cylinder circumference that these applications require.

The irreducible-bridge construction goes back to Kesten and is developed
systematically by Madras and Slade
\cite{Kesten1963,MadrasSlade1993}. The appendix of
\cite{BeatonEtAl2014} already gives its critical normalization and
independent concatenation in the honeycomb mid-edge convention.
Section~\ref{L:step17} derives additional quantitative height and diameter
bounds from our finite boundary estimates. This is a separate renewal
consequence: the proof of Theorem~\ref{L:thm:window} uses finite weighted
counts and recoverable sewing directly.
